\documentclass[11pt,a4paper]{article}

\usepackage[margin=1in]{geometry}
\usepackage{booktabs}
\usepackage{xcolor}
\usepackage{hyperref}
\usepackage{enumitem}
\usepackage{amsmath}
\usepackage{fancyvrb}
\usepackage{titlesec}

\hypersetup{colorlinks=true, linkcolor=blue!50!black, urlcolor=blue!50!black}
\titlespacing*{\section}{0pt}{12pt}{6pt}
\titlespacing*{\subsection}{0pt}{10pt}{4pt}
\setlist{nosep, leftmargin=*}

% Marker for figures that must come from a fresh evaluation run.
% Delete the macro definition body once every one is filled in.
\newcommand{\rerun}[1]{\textcolor{red}{\textbf{[#1]}}}

\title{\vspace{-2em}AI Support Agent for \texttt{@Uber\_Support}\\
\large Hiver SDE Intern Take-Home}
\author{}
\date{}

\begin{document}
\maketitle
\vspace{-3em}

\noindent\rule{\textwidth}{0.4pt}
\smallskip

\noindent\textbf{Dataset:} Customer Support on Twitter (Kaggle, \texttt{thoughtvector/customer-support-on-twitter})\\
\textbf{Brand:} \texttt{@Uber\_Support}, 56{,}307 brand replies in the full file\\
\textbf{Golden set:} 188 hand-labelled examples, \texttt{eval/golden\_set.jsonl}\\
\textbf{Stack:} Qwen3-Embedding-0.6B retrieval, optional Qwen3-Reranker-0.6B, Gemini 2.0 Flash for classification, generation and judging

\noindent\rule{\textwidth}{0.4pt}

\section{Problem framing}

\subsection{What ``good'' means for Uber}

Uber support is not e-commerce support with different nouns. Three things shape what a
good agent looks like here.

\textbf{Safety outranks fluency, and it is not close.} Around 2\% of inbound messages
reference a collision, harassment, a theft or police involvement. A fluent, polite,
wrong reply to any of those is worse than no reply. So the system treats
\texttt{safety\_incident} as a hard stop: it is always escalated, and unlike every other
intent it never gets a draft written at all. Putting a machine-written reply to an
assault report in front of a human to approve is still the machine composing that reply,
and that is not a judgment the system should be making on its own.

\textbf{Twitter is a triage funnel, not a resolution channel.} Reading Uber's own
historical replies makes this obvious. The overwhelming majority redirect to DM or the
in-app help centre, because fares, trip receipts and driver identity cannot be discussed
on a public timeline. So ``good'' is routing correctly and setting expectations, not
resolving the issue in 280 characters. The system is not penalised for failing to solve
anything publicly, because Uber's own agents do not either.

\textbf{No inventing policy.} The generator is grounded in retrieved past resolutions
specifically so it does not improvise refund amounts, ETAs or compensation. A plausible
sounding invented promise is a support liability, not a good reply.

\subsection{What was deliberately not built}

\begin{itemize}
\item \textbf{Multi-turn context.} Only the first brand reply in each thread is used.
Uber threads redirect to DM almost immediately, so later turns are logistics rather than
resolution. This has a real cost, visible in Failure Mode 4.
\item \textbf{Non-English handling.} A character-script filter drops non-Latin-script
messages, mostly Japanese and Hindi. Building multilingual support would have meant
either a translation layer or a multilingual judge, and neither would have been evaluated
properly in the time available. Spanish and other Latin-script languages slip through the
filter, which is a known gap rather than a solved problem.
\item \textbf{Fine-tuning.} Uber's historical agents are inconsistent. Fine-tuning on
their replies would encode that inconsistency as ground truth. Few-shot prompting plus
retrieval keeps the historical data as evidence rather than as a target.
\item \textbf{Banking77.} Offered as an optional secondary source for intent work. Its 77
banking intents do not map onto rideshare traffic, and the brief asks for intents defined
from the data, so borrowing an unrelated taxonomy would have worked against the point.
\end{itemize}

\section{Results against two baselines}

All numbers are computed over the 188-example golden set by \texttt{eval/run\_eval.py}.

\subsection{Intent classification}

\begin{table}[h]
\centering
\begin{tabular}{lccl}
\toprule
\textbf{System} & \textbf{Accuracy} & \textbf{Macro F1} & \textbf{What it is} \\
\midrule
Baseline 1: keyword match & 24.5\% & 0.236 & first substring match wins \\
Baseline 2: TF-IDF + LogReg & 35.1\% & 0.206 & 70/30 split, unigrams and bigrams \\
This system: few-shot LLM & \rerun{rerun} & \rerun{rerun} & structured JSON, safety pre-filter \\
\bottomrule
\end{tabular}
\end{table}

The keyword baseline is deliberately weak and it shows: it cannot place 131 of 188
messages at all and dumps them into \texttt{other}. It also cannot tell ``the app
crashed'' from ``my driver crashed'', which for a rideshare brand is the difference
between a bug report and a collision.

The TF-IDF baseline is the more honest bar. It wins on raw accuracy by learning the head
classes and ignoring the tail, which is exactly what its low macro-F1 (0.206) exposes.
Accuracy alone would flatter it.

\subsection{Escalation and reply quality}

\begin{table}[h]
\centering
\begin{tabular}{lcl}
\toprule
\textbf{Metric} & \textbf{Result} & \textbf{Reading} \\
\midrule
Escalation accuracy & \rerun{rerun} & agreement with human escalate/hold decision \\
Escalation Cohen's $\kappa$ & \rerun{rerun} & the number that actually matters, see \S4 \\
Zero-draft safety routing & \rerun{rerun} & messages that bypassed generation entirely \\
Reply quality, judge mean & \rerun{rerun} & 1 to 5 across four rubric dimensions \\
Judge parse failures & \rerun{rerun} & malformed JSON from the judge \\
\bottomrule
\end{tabular}
\end{table}

\subsection{Is the LLM judge trustworthy}

An LLM judge is worth nothing unless it tracks a human. 30 replies were scored
independently by hand and compared against the judge.

\begin{table}[h]
\centering
\begin{tabular}{lcl}
\toprule
\textbf{Agreement metric} & \textbf{Value} & \textbf{Reading} \\
\midrule
Exact match & 66.7\% & identical score on 20 of 30 \\
Within 1 point & 100.0\% & no catastrophic disagreement \\
Cohen's $\kappa$ & 0.513 & moderate agreement \\
Judge mean vs human mean & 3.27 vs 3.00 & the judge runs lenient \\
\bottomrule
\end{tabular}
\end{table}

The gap between those two means is the useful finding. The judge is systematically
kinder than a human, and consistently so on vague, non-committal replies. Asked to rate
``We've followed up via DM!'' against a customer who wanted a specific answer, the judge
gave 3 and the human gave 2, noting it was evasive. Every reply quality number in this
report should be read with that 0.27 point optimism bias in mind.

\section{Failure analysis}

Five failure modes, with real examples pulled from the evaluation artifacts. Pair IDs
trace back to \texttt{eval/golden\_set.jsonl}.

\subsection{Safety incidents classified as ordinary complaints}

This is the worst failure the system has, so it goes first.

Recall on \texttt{safety\_incident} was 6 of 11. Five genuine safety reports were
classified as something else, and two of those were not escalated at all:

\begin{itemize}
\item \texttt{1662425\_1662424}: ``i've had multiple male drivers record me. is this
allowed? they're all using the same type of camera. very uncomfortable''. Classified as
negative feedback. Not escalated.
\item \texttt{2508288\_2508287}: a rider describing being publicly berated by a driver
over their vision and cognitive ability. Classified as negative feedback. Not escalated.
\item \texttt{1789981\_1789980}: ``my uber driver stole my phone and uber won't help
me''. Classified as a complaint. Escalated, but without the zero-draft guarantee.
\end{itemize}

\textbf{Hypothesis.} The hard keyword pre-filter catches explicit emergency vocabulary
(``police'', ``assault'', ``911'') and the few-shot examples show physical injury. None
of those cover harassment described calmly, covert recording, or theft. The classifier
reads measured language as measured severity. The lesson is that the safety category is
defined by the nature of the event, not the temperature of the writing, and the prompt
does not currently say so.

\subsection{Hallucinated names copied out of retrieved context}

\texttt{330198\_330197}. Customer message in full: ``Dear plz help....''. The system
drafted: ``I'm here to assist you, Boby. Please provide more details about what
assistance is needed.''

There is no Boby. The message carries almost no content, so retrieval returned a weak
nearest neighbour in which a human agent had addressed someone by that name, and the
generator copied it. \textbf{Hypothesis:} with nothing to ground against, the model
treats retrieved text as a template to fill rather than as evidence. Low retrieval
similarity should suppress generation, not proceed with it.

\subsection{Replies that reference links the cleaning step removed}

\texttt{960484\_960482}. Customer: ``hey, so we took a thermus from our driver by mistake
and we don't know how to return it. Help!!''. Draft reply: ``Hi there! Please reach out
to and fill out the 4 boxes at the bottom so we can get in touch.''

``reach out to'' with nothing after it. \textbf{Hypothesis:} \texttt{data\_prep.py}
strips URLs from both sides of every pair. Retrieved examples therefore contain the
grammar of linking without the link, and the generator reproduces the dangling phrasing.
The cleaning step and the generator disagree about what the data means.

\subsection{Single-word follow-ups with no context}

\texttt{1279123\_1279122}. The entire customer message is ``Sent''.

Labelled \texttt{other} by hand because in isolation it genuinely is unclassifiable.
\textbf{Hypothesis:} this is the direct cost of the first-reply-only decision in
\S1.2. ``Sent'' means something specific if the previous turn was ``please DM us your
trip ID'', and nothing at all without it. The system is being asked a question that
cannot be answered from the input it is given, and the golden set label reflects that
honestly rather than pretending a correct answer exists.

\subsection{Safety false positives on ordinary frustration}

The same pre-filter that misses calm harassment reports also over-fires. Six messages
were pushed into \texttt{safety\_incident} that were not: lost keys in a car
(\texttt{order\_delivery}), a Pool driver who did not arrive but charged anyway
(\texttt{billing\_refund}), an account compromise (\texttt{account\_access}).

\textbf{Hypothesis:} substring matching on words like ``police'' and ``accident'' has no
notion of who the sentence is about. These errors are much cheaper than \S3.1, since the
cost is an unnecessary human review rather than a missed incident, and that asymmetry is
deliberate. It is still a precision problem worth naming.

\section{What is misleading about my headline number}

\subsection{The evaluation had a retrieval leak, and the first set of numbers was invalid}

Worth stating plainly because it is the largest thing that went wrong. Golden-set
messages also live in the retrieval index. In the first evaluation run, 176 of 188
predictions retrieved the query's own pair as the top result at similarity exactly 1.0,
which means the generator was shown the ground-truth reply before writing its own.

Two consequences. Reply quality was inflated, because the model was paraphrasing an
answer it had been handed. More subtly, the grounding-similarity branch of the escalation
rule never fired once, because similarity never dropped below the 0.55 threshold. One
third of the escalation policy was dead code and the numbers gave no hint of it.

\texttt{RetrievalIndex.query} now takes \texttt{exclude\_pair\_id} and the pipeline
passes each row's own ID. The affected results are archived unused in
\texttt{eval/stale\_pre\_fix/}. The general lesson is that a similarity of exactly 1.0
appearing in 94\% of rows should have been read as a bug immediately rather than as
strong retrieval.

\subsection{Escalation accuracy is the wrong number to quote}

The golden set splits 94 to 94 on escalate versus hold. On a perfectly balanced set a
coin flip scores 50\% and a system that escalates everything scores 50\%, so raw accuracy
has a floor that looks respectable while meaning nothing. Cohen's $\kappa$ is reported
alongside it for that reason, and $\kappa$ is the number to argue about.

Worth noting that the balance is an artefact of stratified sampling. Real Uber traffic is
nowhere near 50\% escalation-worthy, so neither figure transfers to production volumes
without recalibration.

\subsection{A taxonomy change silently broke the intent numbers}

Partway through, an extra intent (\texttt{feedback\_negative}, splitting venting from
actionable complaints) was added to the classifier. The golden set had already been
labelled without it. The system then predicted a label the ground truth could not contain
55 times, which accounted for 44 of 138 errors and pushed accuracy from roughly 0.40 down
to 0.27.

Nothing in the evaluation flagged this. It looked like a bad model. The split is a real
distinction and worth revisiting, but only together with a relabelling pass, so the intent
was removed to match what the golden set actually encodes.

\subsection{The TF-IDF baseline is trained on 131 examples}

Ten classes, 131 training rows, roughly 13 per class, several classes in single digits.
It is a directional sanity check and nothing more. It should not be read as a ceiling on
what classical methods could do here.

\subsection{$N=188$ is small, and the interval is wide}

A 95\% Wilson interval on a proportion near 0.40 at $N=188$ spans roughly $\pm 7$ points.
Differences smaller than that between systems are noise. \texttt{service\_outage} drew
zero labels in the sample, so any per-class metric for it is computed over an empty
support set and should be ignored rather than interpreted.

\subsection{Nothing here measures faithfulness}

The judge rates whether a reply reads as correct. It does not check the reply against the
retrieved context. A reply that contradicts the historical resolution it was grounded on
can still score well by sounding right, and the ``Boby'' failure in \S3.2 is exactly that
category of error slipping past a quality score.

\section{What one more week would go on}

\begin{enumerate}
\item \textbf{Fix the safety recall problem first.} 6 of 11 is not shippable. Concretely:
add few-shot examples covering harassment, covert recording and theft described in calm
language, and evaluate safety recall as a standalone metric with its own threshold rather
than letting it sit inside macro-F1 where 11 examples cannot move the number.
\item \textbf{Make low retrieval similarity suppress generation.} The threshold currently
only influences escalation. When the best match is weak, the right behaviour is to
produce no draft rather than a confident one built from an unrelated example. That fixes
\S3.2 directly.
\item \textbf{Two turns of thread context.} Enough to resolve ``Sent'' and ``DM done''.
Full dialogue state tracking is not needed for the gain here.
\item \textbf{Tune the escalation thresholds instead of guessing them.} 0.75 and 0.55
were never validated. With a balanced golden set available, a sweep with per-intent
thresholds that weight safety recall above overall accuracy is straightforward.
\item \textbf{Relabel for the venting versus demand split.} It is a genuine distinction
that maps to different handling. It needs the golden set relabelled first, in that order,
which is the lesson from \S4.3.
\end{enumerate}

\vspace{1em}
\noindent\rule{\textwidth}{0.4pt}
\footnotesize
Reproduction: \texttt{README.md}. Decisions and their reasoning: \texttt{decision\_log.md}.
Golden set sampling and labelling method: \texttt{eval/golden\_set\_schema.md}.
All figures regenerate via \texttt{python scripts/report\_numbers.py}.

\end{document}
